\chapter{Introducere}\label{ch:intro}
\pagestyle{fancy}


"Alert fatigue" este fenomenul prin care un operator
expus la un volum mare de alerte, dintre care majoritatea
se dovedesc nerelevante, își pierde treptat capacitatea de a 
reacționa corect la fiecare~\cite{Ndichu2026}. Fenomenul
e documentat extins în securitatea cibernetică: organizațiile
primesc în medie 17.000 de alerte pe săptămână, dintre care
peste 51\% sunt fals pozitive, iar doar 4\% ajung să investigate 
corespunzator~\cite{Hassan2019}

În monitorizarea video, apare acest fenomen, cu efecte similare. 
Sistemele de supraveghere video trimit alerte operatorului 
generate automat de diverși factori precum: declanșare mișcare, 
depășire perimetru, accesarea zonei de interes, detectare prezență umană.
Dacă acese alerte nu sunt atent filtrate ajung sa suprasolicite 
operatorul. Acesta intră în modul de supraviețuire și ajunge, în timp, să
trateze alertele cu indiferență, indiferent de relevanța ei reală.

Acest fenomen de "alert fatigue" poate fi exploatat deliberat 
printr-un atac de tip "cry-wolf". Acest tip de atac descrie scăderea
ratei de răspuns la alarme și alerte pe măsură ce crește proporția de
alarme false primite~\cite{bliss1995reversal}. Un atacator poate genera
în mod repetat alerte false pentru a obosi operatorul, până în punctul în care
o alertă reală este tratată cu aceași superficialitate ca precedentele
false. Consecințele unui astfel de scenariu nu sunt ipotetice: atacul de 
securitate Target din 2013, prin care s-au furat 40 de milioane de înregistrări
de card, s-a produs parțial pentru că personalul nu a reacționat la timp la niște alerte
primite, clasificandu-le greșit drept fals pozitive obișnuite~\cite{Hassan2019}

Când vine vorba de monitorizare video, un posibil atacator poate profita
de oboseala operatorului, declanșând în mod intenționat multiple alarme
de la diverse camere și astfel 



Transpunerea fenomenului din cyber in monitorizare video

Descrierea fenomenului de cry wolf si cum cineva poate profita
de el pentru a obosi operatorul si a efectua un furt (situatia
Dariei)

Problema generala: operatorii care primesc prea multe alerte
incat trec pe modul survival si tind sa ignore
dar asta duce la ratarea unor alerte

Fluxul de lucru: cum vine alerta si cum se gestioneaza

Descriere a experientei mele si statistice obtinute din ce am
discutat cu colegii si am numarat alertele

Net Investigation Velocity: NIV = T-manual \verb|\| T-asistat
(din primul survey citit)

Ideea solutiei:

Descriu cum vreau sa pun un filtru pentru a reduce 
alertele care nu trebuie sa intre (OpenCV, verific daca intr-adevar
e ceva in imagine ce merita sa trimita alerta) si 
dau LM-ului sa analizele a doua oara daca chiar merita 
ca o alerta sa fie trimisa **pe scurt

Domeniul exact:
Tema: verificarea alertelor inainte sa fie trimise la operator
folosind computer vision

Sistemul primeste un clip scurt cu 1 minut inainte de a declansa
alerta si decide daca merita sau nu trimisa spre analiza
operatorului (alerta reala, falsa sau incerta)

Ajuta omul nu il inlocuieste

Nu inlocuieste sistemul initial si nu ia decizii finale de 
interventie (situatii de ce nu ar fi bine sa ia alegerea)

Analog cu gestionarea alertelor SOC, dar aplicata pe video-uri

Structura lucrarii

